\documentclass[11pt,a4paper]{article}
\usepackage[utf8]{inputenc}
\usepackage[T1]{fontenc}
\usepackage[spanish,es-minimal,es-noshorthands]{babel}
\usepackage{lmodern}
\usepackage[sfdefault,lining]{FiraSans}
\usepackage[a4paper,top=2.3cm,bottom=2.2cm,left=1.8cm,right=1.8cm,headheight=1.3cm,headsep=0.4cm,footskip=1cm]{geometry}
\usepackage{amsmath,amssymb}
\usepackage{icomma}
\usepackage{graphicx}
\usepackage[table]{xcolor}
\usepackage{tikz}
\usepackage{fontawesome5}
\usepackage{ccicons}
\usepackage{fancyhdr}
\usepackage{lastpage}
\usepackage{tasks}
\usepackage{needspace}
\usepackage[most]{tcolorbox}
\usepackage{array}
\usepackage[hidelinks]{hyperref}
\AtBeginDocument{\hypersetup{pdftitle={Soluciones de la ficha de la unidad 5: Trigonometría. 4º ESO},pdfsubject={Matemáticas B 4º ESO},pdfkeywords={trigonometría, radianes, seno, coseno, tangente, relaciones fundamentales, triángulos rectángulos, circunferencia goniométrica, reducción al primer cuadrante}}}

\definecolor{navy}{HTML}{1B2748}
\definecolor{accent}{HTML}{E8603A}
\definecolor{teal}{HTML}{1F9E90}
\definecolor{soft}{HTML}{EEF1F8}
\definecolor{muted}{HTML}{5D6685}

\setlength{\parindent}{0pt}
\setlength{\parskip}{3pt}
\settasks{label=\alph*), label-format=\bfseries\color{navy}, label-width=1.3em, item-indent=1.8em, column-sep=0.6em, after-item-skip=5pt, before-skip=3pt, after-skip=4pt}

% ---------- capçalera i peu ----------
\newcommand{\llicencia}{\href{https://creativecommons.org/licenses/by-nc-sa/4.0/deed.es}{\ccbyncsa}\ \ Licencia Creative Commons Atribución-NoComercial-CompartirIgual 4.0 Internacional (CC BY-NC-SA 4.0)}
\pagestyle{fancy}
\fancyhf{}
\fancyhead[L]{\raisebox{-0.35\height}{\includegraphics[height=0.95cm]{logo-gva.png}}}
\fancyhead[C]{\small\color{muted}Soluciones de la ficha · Unidad 5: Trigonometría · 4\textsuperscript{o} ESO B}
\fancyhead[R]{\raisebox{-0.35\height}{\includegraphics[height=0.95cm]{logo-vermellar.png}}}
\fancyfoot[L]{\scriptsize\color{muted}\llicencia}
\fancyfoot[R]{\small\color{muted}\thepage}
\renewcommand{\headrulewidth}{0.4pt}
\renewcommand{\headrule}{\color{muted!40}\hrule height \headrulewidth}
\fancypagestyle{primera}{\fancyhf{}\renewcommand{\headrulewidth}{0pt}
  \fancyfoot[L]{\scriptsize\color{muted}\llicencia}
  \fancyfoot[R]{\small\color{muted}\thepage}}

% ---------- nivells i exercicis ----------
\newcommand{\estrelles}[1]{{\color{accent}\ifcase#1\or\faStar\faStar[regular]\faStar[regular]\or\faStar\faStar\faStar[regular]\or\faStar\faStar\faStar\fi}}
\newcounter{exer}
\newcommand{\ex}[1]{\par\needspace{4\baselineskip}\medskip\refstepcounter{exer}\noindent{\color{navy}\bfseries\theexer.}\ \ {\small\estrelles{#1}}\ \ }

\newcommand{\seccio}[2]{\par\needspace{16\baselineskip}\bigskip
  \begin{tcolorbox}[colback=navy,colframe=navy,coltext=white,arc=2mm,boxrule=0pt,left=3mm,right=3mm,top=1.2mm,bottom=1.2mm]
  {\large\bfseries #1\quad #2}
  \end{tcolorbox}\vspace{-2pt}}

\newcommand{\exf}[1]{\par\medskip\refstepcounter{exer}\noindent{\color{navy}\bfseries\theexer.}\ \ {\small\estrelles{#1}}\ \ }
\newcolumntype{C}[1]{>{\centering\arraybackslash}p{#1}}
% quadre de recordatori (només a la fitxa d'exercicis)
\newcommand{\recorda}[1]{\par\smallskip\begin{tcolorbox}[colback=accent!7,colframe=accent!55,arc=2mm,boxrule=0.5pt,left=2.5mm,right=2.5mm,top=1mm,bottom=1mm]\small{\color{accent}\faLightbulb[regular]}\ \textbf{Recuerda:}\ #1\end{tcolorbox}}
\usetikzlibrary{calc}
\usepackage{enumitem}
\providecommand{\tg}{\operatorname{tg}}
\newcommand{\res}[1]{\boxed{#1}}\newcommand{\rest}[1]{\fbox{#1}}

\def\fmneed{11}% que el títol d'una secció no quede sol al final del full
% =====================================================================
% Capçalera FentMates de les fitxes: un color per a cada curs.
% Els colors són els de la pàgina inicial (LEVELS a assets/js/materials.js);
% el logo ve de _eines/logo/pdf/ (python3 _eines/logo/genera.py).
%
% Ús, al final del preàmbul de la unitat:
%   % =====================================================================
% Capçalera FentMates de les fitxes: un color per a cada curs.
% Els colors són els de la pàgina inicial (LEVELS a assets/js/materials.js);
% el logo ve de _eines/logo/pdf/ (python3 _eines/logo/genera.py).
%
% Ús, al final del preàmbul de la unitat:
%   % =====================================================================
% Capçalera FentMates de les fitxes: un color per a cada curs.
% Els colors són els de la pàgina inicial (LEVELS a assets/js/materials.js);
% el logo ve de _eines/logo/pdf/ (python3 _eines/logo/genera.py).
%
% Ús, al final del preàmbul de la unitat:
%   \input{../comu/fentmates.tex}
%   \fmcurs{1eso}
% i, al començament del document:
%   \capcalera{Matemàtiques}{1r ESO}{Fitxa d'exercicis · Unitat 2}{Divisibilitat}
% Opcional abans de \input: \def\fmneed{14} (espai mínim abans de cada secció).
% =====================================================================
\makeatletter
\def\fm@logodir{../../logo/pdf}
\providecommand{\fmneed}{10}
\def\fm@cursnom{}

% clau del curs, color i color del text damunt del color
\newcommand{\fm@curs}[3]{\@namedef{fm@color@#1}{#2}\@namedef{fm@text@#1}{#3}}
\fm@curs{1eso}{E6B84A}{navy}
\fm@curs{2eso}{3B82F6}{white}
\fm@curs{3eso}{FF7A50}{white}
\fm@curs{4eso-a}{A78BFA}{white}
\fm@curs{4eso-b}{6D28D9}{white}
\fm@curs{1batx-mi}{9B1C1C}{white}
\fm@curs{1batx-ccss}{EF6B6B}{white}
\fm@curs{2batx-mii}{1E6B47}{white}
\fm@curs{2batx-ccss}{2BB3A3}{white}

\newcommand{\fmcurs}[1]{%
  \@ifundefined{fm@color@#1}{\PackageError{fentmates}{Curs desconegut: #1}{}}{}%
  \edef\fm@tmp{\noexpand\definecolor{curs}{HTML}{\@nameuse{fm@color@#1}}%
               \noexpand\colorlet{curstext}{\@nameuse{fm@text@#1}}}\fm@tmp
  \colorlet{cursfosc}{curs!72!black}%
  \def\fm@logo{\fm@logodir/vertical-#1.pdf}%
  \def\fm@marca{\fm@logodir/marca-#1.pdf}}

% ---------- capçalera de la primera pàgina ----------
% #1 matèria (petit), #2 curs, #3 tipus de document i unitat, #4 títol
\newcommand{\capcalera}[4]{%
  \gdef\fm@cursnom{#2}%
  \par\noindent
  \begin{tcolorbox}[colback=curs,colframe=curs,boxrule=0pt,arc=2.2mm,width=3.5cm,height=2.15cm,
      valign=center,left=3mm,right=2mm,top=1mm,bottom=1mm,nobeforeafter,box align=center]
    \color{curstext}{\footnotesize #1}\par\vspace{1pt}{\fontsize{19}{21}\selectfont\bfseries #2\par}%
  \end{tcolorbox}\hspace{5mm}%
  \begin{minipage}[c]{\dimexpr\linewidth-3.5cm-5mm-2.9cm\relax}
    {\color{cursfosc}\small\bfseries\MakeUppercase{#3}}\par\vspace{2pt}
    {\color{navy}\fontsize{20}{23}\selectfont\bfseries #4\par}
  \end{minipage}\hfill
  \begin{minipage}[c]{2.6cm}\raggedleft\includegraphics[width=2.6cm]{\fm@logo}\end{minipage}\par
  \vspace{2pt}}

% ---------- barres de secció del color del curs ----------
\renewcommand{\seccio}[2]{\par\needspace{\fmneed\baselineskip}\bigskip
  \begin{tcolorbox}[colback=curs,colframe=curs,coltext=curstext,arc=2mm,boxrule=0pt,left=3mm,right=3mm,top=1.2mm,bottom=1.2mm]
  {\large\bfseries #1\quad #2}
  \end{tcolorbox}\vspace{-2pt}}

% ---------- estrelles de dificultat del color del curs (a Batxillerat no se'n posen) ----------
\@ifundefined{estrelles}{}{%
  \renewcommand{\estrelles}[1]{{\color{curs}\ifcase#1\or\faStar\faStar[regular]\faStar[regular]\or\faStar\faStar\faStar[regular]\or\faStar\faStar\faStar\fi}}}

% ---------- capçalera i peu de les pàgines ----------
% primera pàgina: logos institucionals al peu
\fancypagestyle{primera}{\fancyhf{}\renewcommand{\headrulewidth}{0pt}
  \fancyfoot[L]{\raisebox{-0.32\height}{\includegraphics[height=0.78cm]{logo-gva.png}}\hspace{5mm}%
    \raisebox{-0.32\height}{\includegraphics[height=0.78cm]{logo-vermellar.png}}}
  \fancyfoot[R]{\scriptsize\color{muted}\href{https://creativecommons.org/licenses/by-nc-sa/4.0/deed.ca}{\ccbyncsa}\ \ CC BY-NC-SA 4.0 · fentmates.org\qquad\small\thepage}}
% la resta: el símbol a l'esquerra i el curs a la dreta, amb el títol al centre com abans
\fancyhead[L]{\raisebox{-0.3\height}{\includegraphics[height=0.62cm]{\fm@marca}}}
\fancyhead[R]{\small\bfseries\color{cursfosc}\fm@cursnom}
\renewcommand{\headrule}{\color{curs!45}\hrule height \headrulewidth}
\makeatother

%   \fmcurs{1eso}
% i, al començament del document:
%   \capcalera{Matemàtiques}{1r ESO}{Fitxa d'exercicis · Unitat 2}{Divisibilitat}
% Opcional abans de \input: \def\fmneed{14} (espai mínim abans de cada secció).
% =====================================================================
\makeatletter
\def\fm@logodir{../../logo/pdf}
\providecommand{\fmneed}{10}
\def\fm@cursnom{}

% clau del curs, color i color del text damunt del color
\newcommand{\fm@curs}[3]{\@namedef{fm@color@#1}{#2}\@namedef{fm@text@#1}{#3}}
\fm@curs{1eso}{E6B84A}{navy}
\fm@curs{2eso}{3B82F6}{white}
\fm@curs{3eso}{FF7A50}{white}
\fm@curs{4eso-a}{A78BFA}{white}
\fm@curs{4eso-b}{6D28D9}{white}
\fm@curs{1batx-mi}{9B1C1C}{white}
\fm@curs{1batx-ccss}{EF6B6B}{white}
\fm@curs{2batx-mii}{1E6B47}{white}
\fm@curs{2batx-ccss}{2BB3A3}{white}

\newcommand{\fmcurs}[1]{%
  \@ifundefined{fm@color@#1}{\PackageError{fentmates}{Curs desconegut: #1}{}}{}%
  \edef\fm@tmp{\noexpand\definecolor{curs}{HTML}{\@nameuse{fm@color@#1}}%
               \noexpand\colorlet{curstext}{\@nameuse{fm@text@#1}}}\fm@tmp
  \colorlet{cursfosc}{curs!72!black}%
  \def\fm@logo{\fm@logodir/vertical-#1.pdf}%
  \def\fm@marca{\fm@logodir/marca-#1.pdf}}

% ---------- capçalera de la primera pàgina ----------
% #1 matèria (petit), #2 curs, #3 tipus de document i unitat, #4 títol
\newcommand{\capcalera}[4]{%
  \gdef\fm@cursnom{#2}%
  \par\noindent
  \begin{tcolorbox}[colback=curs,colframe=curs,boxrule=0pt,arc=2.2mm,width=3.5cm,height=2.15cm,
      valign=center,left=3mm,right=2mm,top=1mm,bottom=1mm,nobeforeafter,box align=center]
    \color{curstext}{\footnotesize #1}\par\vspace{1pt}{\fontsize{19}{21}\selectfont\bfseries #2\par}%
  \end{tcolorbox}\hspace{5mm}%
  \begin{minipage}[c]{\dimexpr\linewidth-3.5cm-5mm-2.9cm\relax}
    {\color{cursfosc}\small\bfseries\MakeUppercase{#3}}\par\vspace{2pt}
    {\color{navy}\fontsize{20}{23}\selectfont\bfseries #4\par}
  \end{minipage}\hfill
  \begin{minipage}[c]{2.6cm}\raggedleft\includegraphics[width=2.6cm]{\fm@logo}\end{minipage}\par
  \vspace{2pt}}

% ---------- barres de secció del color del curs ----------
\renewcommand{\seccio}[2]{\par\needspace{\fmneed\baselineskip}\bigskip
  \begin{tcolorbox}[colback=curs,colframe=curs,coltext=curstext,arc=2mm,boxrule=0pt,left=3mm,right=3mm,top=1.2mm,bottom=1.2mm]
  {\large\bfseries #1\quad #2}
  \end{tcolorbox}\vspace{-2pt}}

% ---------- estrelles de dificultat del color del curs (a Batxillerat no se'n posen) ----------
\@ifundefined{estrelles}{}{%
  \renewcommand{\estrelles}[1]{{\color{curs}\ifcase#1\or\faStar\faStar[regular]\faStar[regular]\or\faStar\faStar\faStar[regular]\or\faStar\faStar\faStar\fi}}}

% ---------- capçalera i peu de les pàgines ----------
% primera pàgina: logos institucionals al peu
\fancypagestyle{primera}{\fancyhf{}\renewcommand{\headrulewidth}{0pt}
  \fancyfoot[L]{\raisebox{-0.32\height}{\includegraphics[height=0.78cm]{logo-gva.png}}\hspace{5mm}%
    \raisebox{-0.32\height}{\includegraphics[height=0.78cm]{logo-vermellar.png}}}
  \fancyfoot[R]{\scriptsize\color{muted}\href{https://creativecommons.org/licenses/by-nc-sa/4.0/deed.ca}{\ccbyncsa}\ \ CC BY-NC-SA 4.0 · fentmates.org\qquad\small\thepage}}
% la resta: el símbol a l'esquerra i el curs a la dreta, amb el títol al centre com abans
\fancyhead[L]{\raisebox{-0.3\height}{\includegraphics[height=0.62cm]{\fm@marca}}}
\fancyhead[R]{\small\bfseries\color{cursfosc}\fm@cursnom}
\renewcommand{\headrule}{\color{curs!45}\hrule height \headrulewidth}
\makeatother

%   \fmcurs{1eso}
% i, al començament del document:
%   \capcalera{Matemàtiques}{1r ESO}{Fitxa d'exercicis · Unitat 2}{Divisibilitat}
% Opcional abans de \input: \def\fmneed{14} (espai mínim abans de cada secció).
% =====================================================================
\makeatletter
\def\fm@logodir{../../logo/pdf}
\providecommand{\fmneed}{10}
\def\fm@cursnom{}

% clau del curs, color i color del text damunt del color
\newcommand{\fm@curs}[3]{\@namedef{fm@color@#1}{#2}\@namedef{fm@text@#1}{#3}}
\fm@curs{1eso}{E6B84A}{navy}
\fm@curs{2eso}{3B82F6}{white}
\fm@curs{3eso}{FF7A50}{white}
\fm@curs{4eso-a}{A78BFA}{white}
\fm@curs{4eso-b}{6D28D9}{white}
\fm@curs{1batx-mi}{9B1C1C}{white}
\fm@curs{1batx-ccss}{EF6B6B}{white}
\fm@curs{2batx-mii}{1E6B47}{white}
\fm@curs{2batx-ccss}{2BB3A3}{white}

\newcommand{\fmcurs}[1]{%
  \@ifundefined{fm@color@#1}{\PackageError{fentmates}{Curs desconegut: #1}{}}{}%
  \edef\fm@tmp{\noexpand\definecolor{curs}{HTML}{\@nameuse{fm@color@#1}}%
               \noexpand\colorlet{curstext}{\@nameuse{fm@text@#1}}}\fm@tmp
  \colorlet{cursfosc}{curs!72!black}%
  \def\fm@logo{\fm@logodir/vertical-#1.pdf}%
  \def\fm@marca{\fm@logodir/marca-#1.pdf}}

% ---------- capçalera de la primera pàgina ----------
% #1 matèria (petit), #2 curs, #3 tipus de document i unitat, #4 títol
\newcommand{\capcalera}[4]{%
  \gdef\fm@cursnom{#2}%
  \par\noindent
  \begin{tcolorbox}[colback=curs,colframe=curs,boxrule=0pt,arc=2.2mm,width=3.5cm,height=2.15cm,
      valign=center,left=3mm,right=2mm,top=1mm,bottom=1mm,nobeforeafter,box align=center]
    \color{curstext}{\footnotesize #1}\par\vspace{1pt}{\fontsize{19}{21}\selectfont\bfseries #2\par}%
  \end{tcolorbox}\hspace{5mm}%
  \begin{minipage}[c]{\dimexpr\linewidth-3.5cm-5mm-2.9cm\relax}
    {\color{cursfosc}\small\bfseries\MakeUppercase{#3}}\par\vspace{2pt}
    {\color{navy}\fontsize{20}{23}\selectfont\bfseries #4\par}
  \end{minipage}\hfill
  \begin{minipage}[c]{2.6cm}\raggedleft\includegraphics[width=2.6cm]{\fm@logo}\end{minipage}\par
  \vspace{2pt}}

% ---------- barres de secció del color del curs ----------
\renewcommand{\seccio}[2]{\par\needspace{\fmneed\baselineskip}\bigskip
  \begin{tcolorbox}[colback=curs,colframe=curs,coltext=curstext,arc=2mm,boxrule=0pt,left=3mm,right=3mm,top=1.2mm,bottom=1.2mm]
  {\large\bfseries #1\quad #2}
  \end{tcolorbox}\vspace{-2pt}}

% ---------- estrelles de dificultat del color del curs (a Batxillerat no se'n posen) ----------
\@ifundefined{estrelles}{}{%
  \renewcommand{\estrelles}[1]{{\color{curs}\ifcase#1\or\faStar\faStar[regular]\faStar[regular]\or\faStar\faStar\faStar[regular]\or\faStar\faStar\faStar\fi}}}

% ---------- capçalera i peu de les pàgines ----------
% primera pàgina: logos institucionals al peu
\fancypagestyle{primera}{\fancyhf{}\renewcommand{\headrulewidth}{0pt}
  \fancyfoot[L]{\raisebox{-0.32\height}{\includegraphics[height=0.78cm]{logo-gva.png}}\hspace{5mm}%
    \raisebox{-0.32\height}{\includegraphics[height=0.78cm]{logo-vermellar.png}}}
  \fancyfoot[R]{\scriptsize\color{muted}\href{https://creativecommons.org/licenses/by-nc-sa/4.0/deed.ca}{\ccbyncsa}\ \ CC BY-NC-SA 4.0 · fentmates.org\qquad\small\thepage}}
% la resta: el símbol a l'esquerra i el curs a la dreta, amb el títol al centre com abans
\fancyhead[L]{\raisebox{-0.3\height}{\includegraphics[height=0.62cm]{\fm@marca}}}
\fancyhead[R]{\small\bfseries\color{cursfosc}\fm@cursnom}
\renewcommand{\headrule}{\color{curs!45}\hrule height \headrulewidth}
\makeatother

\fmcurs{4eso-b}

\begin{document}
\thispagestyle{primera}

\capcalera{Matemáticas B}{4\textsuperscript{o} ESO}{Soluciones de la ficha · Unidad 5}{\raggedright Trigonometría}
\vspace{8pt}

\begin{tcolorbox}[colback=soft,colframe=soft,arc=2mm,left=3mm,right=3mm,top=1.5mm,bottom=1.5mm]
\small Intenta resolver cada ejercicio antes de mirar su solución. Los ángulos se han redondeado al segundo y las longitudes, a la centésima (calculando siempre con los valores sin redondear). Puedes comprobar muchos resultados con las herramientas interactivas de los apuntes de la unidad.
\end{tcolorbox}

% =====================================================
\seccio{1}{Medida de ángulos: grados y radianes}

\exf{1} \begin{tasks}(6)
  \task $\res{\frac{\pi}{4}}$
  \task $\res{\frac{\pi}{6}}$
  \task $\res{\frac{2\pi}{3}}$
  \task $\res{\frac{5\pi}{6}}$
  \task $\res{3\pi}$
  \task $\res{6\pi}$
  \task $\res{10\pi}$
  \task $\res{\frac{4\pi}{3}}$
  \task $\res{\frac{31\pi}{18}}$
  \task $\res{\frac{7\pi}{18}}$
  \task $\res{\frac{\pi}{15}}$
  \task $\res{4\pi}$
\end{tasks}

\ex{1} \begin{tasks}(6)
  \task $\res{180^\circ}$
  \task $\res{720^\circ}$
  \task $\res{1260^\circ}$
  \task $\res{30^\circ}$
  \task $\res{135^\circ}$
  \task $\res{270^\circ}$
  \task $\res{300^\circ}$
  \task $\res{90^\circ}$
  \task $\res{240^\circ}$
  \task $\res{216^\circ}$
  \task $\res{72^\circ}$
  \task $\res{20^\circ}$
\end{tasks}

\ex{1} \begin{center}\renewcommand{\arraystretch}{1.6}\begin{tabular}{|l|*{10}{C{1.05cm}|}}\hline\textbf{Grados} & $30^\circ$ & $\res{45^\circ}$ & $60^\circ$ & $\res{120^\circ}$ & $135^\circ$ & $\res{150^\circ}$ & $210^\circ$ & $\res{225^\circ}$ & $300^\circ$ & $\res{330^\circ}$\\\hline\textbf{Radianes} & $\res{\frac{\pi}{6}}$ & $\frac{\pi}{4}$ & $\res{\frac{\pi}{3}}$ & $\frac{2\pi}{3}$ & $\res{\frac{3\pi}{4}}$ & $\frac{5\pi}{6}$ & $\res{\frac{7\pi}{6}}$ & $\frac{5\pi}{4}$ & $\res{\frac{5\pi}{3}}$ & $\frac{11\pi}{6}$\\\hline\end{tabular}\end{center}

\ex{1} \begin{tasks}(4)
  \task $\res{90^\circ=\frac{\pi}{2}}$
  \task $\res{270^\circ=\frac{3\pi}{2}}$
  \task $\res{60^\circ=\frac{\pi}{3}}$
  \task $\res{120^\circ=\frac{2\pi}{3}}$
\end{tasks}

\ex{2} $\widehat{AOB}=\frac{360^\circ}{6}=\res{\frac{\pi}{3}}$; \ ángulo interior: $\widehat{ABC}=120^\circ=\res{\frac{2\pi}{3}}$; \ el triángulo $ABC$ es isósceles, y $\widehat{BAC}=\frac{180^\circ-120^\circ}{2}=30^\circ=\res{\frac{\pi}{6}}$; \ $AD$ es un diámetro y el ángulo inscrito $\widehat{ACD}$ es recto: $\res{\frac{\pi}{2}}$.

\ex{2} \begin{tasks}(3)
  \task $\res{57^\circ17'45''}$
  \task $\res{42^\circ58'19''}$
  \task $\res{143^\circ14'22''}$
  \task $\res{\frac{\pi}{8}}$
  \task $\res{\frac{7\pi}{24}}$
  \task $\res{\frac{7\pi}{8}}$
\end{tasks}

\ex{3} Horas: $\frac{360^\circ}{12}=\res{30^\circ=\frac{\pi}{6}}$. Minutos: $\frac{20}{60}\cdot360^\circ=\res{120^\circ=\frac{2\pi}{3}}$. A las 4: $4\cdot30^\circ=\res{120^\circ=\frac{2\pi}{3}}$.


% =====================================================
\seccio{2}{Razones trigonométricas de un ángulo agudo}

\exf{1} \begin{enumerate}[label=\alph*),leftmargin=*,itemsep=2pt]\item Pitágoras: hipotenusa $34$ cm. \ $\sin\alpha=\res{\frac{8}{17}}$, $\cos\alpha=\res{\frac{15}{17}}$, $\tg\alpha=\res{\frac{8}{15}}$; \ $\sin\beta=\res{\frac{15}{17}}$, $\cos\beta=\res{\frac{8}{17}}$, $\tg\beta=\res{\frac{15}{8}}$\item Pitágoras: cateto contiguo $8$ cm. \ $\sin\alpha=\res{\frac{3}{5}}$, $\cos\alpha=\res{\frac{4}{5}}$, $\tg\alpha=\res{\frac{3}{4}}$; \ $\sin\beta=\res{\frac{4}{5}}$, $\cos\beta=\res{\frac{3}{5}}$, $\tg\beta=\res{\frac{4}{3}}$\item Pitágoras: hipotenusa $26$ cm. \ $\sin\alpha=\res{\frac{12}{13}}$, $\cos\alpha=\res{\frac{5}{13}}$, $\tg\alpha=\res{\frac{12}{5}}$; \ $\sin\beta=\res{\frac{5}{13}}$, $\cos\beta=\res{\frac{12}{13}}$, $\tg\beta=\res{\frac{5}{12}}$\item Pitágoras: cateto opuesto $24$ cm. \ $\sin\alpha=\res{\frac{24}{25}}$, $\cos\alpha=\res{\frac{7}{25}}$, $\tg\alpha=\res{\frac{24}{7}}$; \ $\sin\beta=\res{\frac{7}{25}}$, $\cos\beta=\res{\frac{24}{25}}$, $\tg\beta=\res{\frac{7}{24}}$\item Pitágoras: cateto opuesto $21$ cm. \ $\sin\alpha=\res{\frac{21}{29}}$, $\cos\alpha=\res{\frac{20}{29}}$, $\tg\alpha=\res{\frac{21}{20}}$; \ $\sin\beta=\res{\frac{20}{29}}$, $\cos\beta=\res{\frac{21}{29}}$, $\tg\beta=\res{\frac{20}{21}}$\end{enumerate}

\ex{1} Pitágoras: $\overline{AD}=\sqrt{15^2-12^2}=9$ cm y $\overline{BC}=\sqrt{12^2+16^2}=20$ cm.\begin{enumerate}[label=\alph*),leftmargin=*,itemsep=2pt]\item $\sin\hat A=\frac{12}{15}=\res{\frac45}$, \ $\cos\hat A=\frac{9}{15}=\res{\frac35}$, \ $\tg\hat A=\frac{12}{9}=\res{\frac43}$\item $\sin\hat C=\frac{12}{20}=\res{\frac35}$, \ $\cos\hat C=\frac{16}{20}=\res{\frac45}$, \ $\tg\hat C=\frac{12}{16}=\res{\frac34}$\item $\sin\widehat{ABD}=\frac{9}{15}=\res{\frac35}$, \ $\cos\widehat{ABD}=\frac{12}{15}=\res{\frac45}$, \ $\tg\widehat{ABD}=\frac{9}{12}=\res{\frac34}$\item $\sin\widehat{CBD}=\frac{16}{20}=\res{\frac45}$, \ $\cos\widehat{CBD}=\frac{12}{20}=\res{\frac35}$, \ $\tg\widehat{CBD}=\frac{16}{12}=\res{\frac43}$\end{enumerate}

\ex{1} \begin{center}\renewcommand{\arraystretch}{1.45}\begin{tabular}{|c|*{4}{C{2.3cm}|}}\hline $\alpha$ & $15^\circ$ & $55^\circ20'$ & $72^\circ25'40''$ & $85{,}5^\circ$\\\hline
$\sin\alpha$ & $\res{0,2588}$ & $\res{0,8225}$ & $\res{0,9533}$ & $\res{0,9969}$\\\hline $\cos\alpha$ & $\res{0,9659}$ & $\res{0,5688}$ & $\res{0,3019}$ & $\res{0,0785}$\\\hline $\tg\alpha$ & $\res{0,2679}$ & $\res{1,4460}$ & $\res{3,1577}$ & $\res{12,7062}$\\\hline\end{tabular}\end{center}

\ex{1} \begin{tasks}(3)
  \task $\alpha=\res{60^\circ}$
  \task $\alpha=\approx\res{70^\circ31'44''}$
  \task $\alpha=\approx\res{68^\circ11'55''}$
  \task $\alpha=\approx\res{14^\circ28'39''}$
  \task $\alpha=\approx\res{36^\circ52'12''}$
\end{tasks}

\ex{2} Medio triángulo equilátero de lado 4: catetos 2 y $\sqrt{16-4}=2\sqrt3$. Medio triángulo de lado 6: catetos 3 y $\sqrt{36-9}=3\sqrt3$. En los dos: $\sin60^\circ=\frac{2\sqrt3}{4}=\frac{3\sqrt3}{6}=\res{\frac{\sqrt3}{2}}$, \ $\cos60^\circ=\frac24=\frac36=\res{\frac12}$, \ $\tg60^\circ=\frac{2\sqrt3}{2}=\frac{3\sqrt3}{3}=\res{\sqrt3}$: \rest{no dependen del lado}.

\ex{2} \begin{enumerate}[label=\alph*),leftmargin=*,itemsep=2pt]\item Cateto opuesto: $5\cdot\frac65=\res{6\ \text{cm}}$; \ hipotenusa: $\sqrt{25+36}=\res{\sqrt{61}\approx7,81\ \text{cm}}$.\item Hipotenusa: $\frac{24}{12/13}=\res{26\ \text{cm}}$; \ cateto: $\sqrt{26^2-24^2}=\res{10\ \text{cm}}$; \ $\cos=\res{\frac{5}{13}}$, \ $\tg=\res{\frac{12}{5}}$.\item Cateto contiguo: $\frac{15}{3/4}=\res{20\ \text{cm}}$; \ hipotenusa: $\sqrt{15^2+20^2}=\res{25\ \text{cm}}$; \ $\sin=\res{\frac35}$, \ $\cos=\res{\frac45}$.\item Los dos catetos son iguales: $x^2+x^2=4\Rightarrow x=\res{\sqrt2\ \text{cm}}$ cada uno.\end{enumerate}

\ex{3} Los catetos son $7x$ y $12x$: $49x^2+144x^2=100\Rightarrow x=\frac{10}{\sqrt{193}}$. Catetos: $\res{\frac{70\sqrt{193}}{193}\approx5,04\ \text{cm}}$ y $\res{\frac{120\sqrt{193}}{193}\approx8,64\ \text{cm}}$. \ $\sin\alpha=\res{\frac{7\sqrt{193}}{193}}$, \ $\cos\alpha=\res{\frac{12\sqrt{193}}{193}}$.


% =====================================================
\seccio{3}{Razones trigonométricas de 30°, 45° y 60°}

\exf{1} \begin{center}\renewcommand{\arraystretch}{1.6}\begin{tabular}{|l|*{5}{C{1.6cm}|}}\hline
& $0^\circ$ & $30^\circ$ & $45^\circ$ & $60^\circ$ & $90^\circ$\\\hline
$\sin$ & $0$ & $\frac12$ & $\frac{\sqrt2}{2}$ & $\frac{\sqrt3}{2}$ & $1$\\\hline $\cos$ & $1$ & $\frac{\sqrt3}{2}$ & $\frac{\sqrt2}{2}$ & $\frac12$ & $0$\\\hline $\tg$ & $0$ & $\frac{\sqrt3}{3}$ & $1$ & $\sqrt3$ & $\nexists$\\\hline\end{tabular}\end{center}

\ex{1} \begin{enumerate}[label=\alph*),leftmargin=*,itemsep=2pt]\item Hipotenusa: $\res{\frac{10\sqrt3}{3}\ \text{cm}}$; \ el cateto contiguo: $\res{\frac{5\sqrt3}{3}\ \text{cm}}$.\item Hipotenusa: $\res{4\ \text{cm}}$; \ el cateto contiguo: $\res{2\sqrt3\ \text{cm}}$.\item Hipotenusa: $\res{7\sqrt2\ \text{cm}}$; \ el cateto opuesto: $\res{7\ \text{cm}}$.\item Hipotenusa: $\res{12\ \text{cm}}$; \ el cateto opuesto: $\res{6\sqrt3\ \text{cm}}$.\item Hipotenusa: $\res{12\sqrt2\ \text{cm}}$; \ el cateto contiguo: $\res{12\ \text{cm}}$.\item Cateto opuesto: $30\sin30^\circ=\res{15\ \text{cm}}$; \ contiguo: $30\cos30^\circ=\res{15\sqrt3\ \text{cm}}$.\end{enumerate}

\ex{1} \begin{tasks}(3)
  \task $\res{1}$
  \task $\res{0}$
  \task $\res{\frac12}$
  \task $\res{1}$
  \task $\res{2}$
  \task $\res{\frac12}$
\end{tasks}

\ex{2} Altura: $8\sin60^\circ=8\cdot\frac{\sqrt3}{2}=\res{4\sqrt3\approx6,93\ \text{cm}}$. \ Diagonal: $\cos45^\circ=\frac5d\Rightarrow d=\frac{5}{\sqrt2/2}=\res{5\sqrt2\approx7,07\ \text{cm}}$.

\ex{2} La apotema es la altura de un triángulo equilátero de lado 4: $\text{ap}=4\sin60^\circ=\res{2\sqrt3\ \text{cm}}$. \ Área: $6\cdot\frac{4\cdot2\sqrt3}{2}=\res{24\sqrt3\approx41,57\ \text{cm}^2}$.

\ex{3} Antena $PB$: $\overline{AB}=\frac{100}{\sin60^\circ}=\res{\frac{200\sqrt3}{3}\ \text{m}}$, $\overline{AP}=\frac{100}{\tg60^\circ}=\frac{100\sqrt3}{3}$; \ $\overline{BC}=\frac{100}{\sin30^\circ}=\res{200\ \text{m}}$, $\overline{PC}=\frac{100}{\tg30^\circ}=100\sqrt3$. \ Antena $QD$ (los ángulos son los de $D$): $\overline{DC}=\frac{75}{\cos45^\circ}=\res{75\sqrt2\ \text{m}}$, $\overline{CQ}=75\tg45^\circ=75$; \ $\overline{DE}=\frac{75}{\cos30^\circ}=\res{50\sqrt3\ \text{m}}$, $\overline{QE}=75\tg30^\circ=25\sqrt3$. \ $\overline{AE}=\frac{100\sqrt3}{3}+100\sqrt3+75+25\sqrt3=\res{75+\frac{475\sqrt3}{3}\approx349,24\ \text{m}}$.


% =====================================================
\seccio{4}{Relaciones fundamentales}

\exf{1} \begin{tasks}(2)
  \task $\cos\alpha=\res{\frac{24}{25}}$, \ $\tg\alpha=\res{\frac{7}{24}}$
  \task $\cos\alpha=\res{\frac{35}{37}}$, \ $\tg\alpha=\res{\frac{12}{35}}$
  \task $\cos\alpha=\res{\frac{60}{61}}$, \ $\tg\alpha=\res{\frac{11}{60}}$
  \task $\sin\alpha=\res{\frac{21}{29}}$, \ $\tg\alpha=\res{\frac{21}{20}}$
  \task $\sin\alpha=\res{\frac{5}{13}}$, \ $\tg\alpha=\res{\frac{5}{12}}$
  \task $\cos\alpha=\res{\frac{\sqrt{55}}{8}}$, \ $\tg\alpha=\res{\frac{3 \sqrt{55}}{55}}$
  \task $\sin\alpha=\res{\frac{2 \sqrt{2}}{3}}$, \ $\tg\alpha=\res{2 \sqrt{2}}$
  \task $\cos\alpha=\res{\frac{7 \sqrt{19}}{34}}$, \ $\tg\alpha=\res{\frac{15 \sqrt{19}}{133}}$
  \task $\sin\alpha=\res{\frac{\sqrt{21}}{5}}$, \ $\tg\alpha=\res{\frac{\sqrt{21}}{2}}$
  \task $\cos\alpha=\res{\frac{\sqrt{33}}{7}}$, \ $\tg\alpha=\res{\frac{4 \sqrt{33}}{33}}$
\end{tasks}

\ex{2} \begin{tasks}(2)
  \task $\sin\alpha=\res{\frac{5}{13}}$, \ $\cos\alpha=\res{\frac{12}{13}}$
  \task $\sin\alpha=\res{\frac{3}{5}}$, \ $\cos\alpha=\res{\frac{4}{5}}$
  \task $\sin\alpha=\res{\frac{5 \sqrt{34}}{34}}$, \ $\cos\alpha=\res{\frac{3 \sqrt{34}}{34}}$
  \task $\sin\alpha=\res{\frac{5 \sqrt{194}}{194}}$, \ $\cos\alpha=\res{\frac{13 \sqrt{194}}{194}}$
  \task $\sin\alpha=\res{\frac{5 \sqrt{89}}{89}}$, \ $\cos\alpha=\res{\frac{8 \sqrt{89}}{89}}$
  \task $\sin\alpha=\res{\frac{3 \sqrt{10}}{10}}$, \ $\cos\alpha=\res{\frac{\sqrt{10}}{10}}$
  \task $\sin\alpha=\res{\frac{\sqrt{2}}{2}}$, \ $\cos\alpha=\res{\frac{\sqrt{2}}{2}}$
\end{tasks}

\ex{1} \begin{tasks}(2)
  \task $\sin\alpha\approx\res{0,9907}$, \ $\tg\alpha\approx\res{7,3010}$
  \task $\sin\alpha\approx\res{0,3027}$, \ $\tg\alpha\approx\res{0,3175}$
  \task $\sin\alpha\approx\res{0,9581}$, \ $\tg\alpha\approx\res{3,3454}$
  \task $\cos\alpha\approx\res{0,9832}$, \ $\tg\alpha\approx\res{0,1858}$
  \task $\cos\alpha\approx\res{0,2992}$, \ $\tg\alpha\approx\res{3,1895}$
  \task $\cos\alpha\approx\res{0,8915}$, \ $\tg\alpha\approx\res{0,5083}$
\end{tasks}

\ex{2} \begin{tasks}(1)
  \task $\frac{49}{625}+\frac{121}{625}=\frac{170}{625}\neq1$: \rest{no hay ninguno}
  \task $\frac{49}{169}+\frac{144}{169}=\frac{193}{169}\neq1$: \rest{no hay ninguno}
  \task $\frac{5}{25}+\frac{20}{25}=1$: \rest{sí que hay}
  \task $\frac{9}{25}+\frac{16}{25}=1$: \rest{sí que hay}
  \task $\cos\alpha=\sqrt{1-\tfrac{49}{121}}=\tfrac{6\sqrt2}{11}\ \text{y}\ \tfrac{\sin\alpha}{\cos\alpha}=\tfrac{7}{6\sqrt2}\neq\tfrac{6}{11}$: \rest{no hay ninguno}
  \task $1^2+1=2\neq\tfrac{169}{25}=\tfrac{1}{\cos^2\alpha}$: \rest{no hay ninguno}
  \task $\tfrac39+1=\tfrac43=\tfrac{1}{3/4}$: \rest{sí que hay}
  \task $\tfrac{49}{25}+1=\tfrac{74}{25}\neq\tfrac{25}{9}$: \rest{no hay ninguno}
\end{tasks}

\ex{2} \begin{tasks}(1)
  \task $0,4663^2+1\approx1,2174$, \ $\frac{1}{0,1736^2}\approx33,1818$: \rest{no}
  \task $0,3640^2+1\approx1,1325$, \ $\frac{1}{0,9397^2}\approx1,1325$: \rest{sí}
  \task $0,1485^2+0,2588^2\approx0,0890$: \rest{no}
  \task $0,9325^2+1\approx1,8696$, \ $\frac{1}{0,7313^2}\approx1,8699$: \rest{sí}
  \task $0,2754^2+0,6691^2\approx0,5235$: \rest{no}
\end{tasks}

\ex{3} \begin{tasks}(1)
  \task Dividimos $\sin^2\alpha+\cos^2\alpha=1$ entre $\sin^2\alpha$: $1+\dfrac{\cos^2\alpha}{\sin^2\alpha}=\dfrac{1}{\sin^2\alpha}$, y $\dfrac{\cos^2\alpha}{\sin^2\alpha}=\dfrac{1}{\tg^2\alpha}$.
  \task $(1+\tg^2\alpha)\cos^2\alpha=\cos^2\alpha+\dfrac{\sin^2\alpha}{\cos^2\alpha}\cos^2\alpha=\cos^2\alpha+\sin^2\alpha=1$.
  \task $\sin\alpha\cos\alpha\dfrac{\sin\alpha}{\cos\alpha}+\cos^2\alpha=\sin^2\alpha+\cos^2\alpha=1$.
\end{tasks}


% =====================================================
\seccio{5}{Resolución de triángulos rectángulos}

\exf{1} \begin{tasks}(1)
  \task $\hat C=\res{60^\circ}$, \ $a=\frac{4}{\sin30^\circ}=\res{8\ \text{m}}$, \ $c=\frac{4}{\tg30^\circ}=\res{4\sqrt3\approx6,93\ \text{m}}$
  \task $c=\sqrt{7^2-5^2}=\res{2\sqrt6\approx4,9\ \text{m}}$, \ $\sin\hat B=\frac57\Rightarrow\hat B\approx\res{45^\circ35'5''}$, \ $\hat C\approx\res{44^\circ24'55''}$
  \task $\hat C=\res{50^\circ}$, \ $b=10\sin40^\circ\approx\res{6,43\ \text{m}}$, \ $c=10\cos40^\circ\approx\res{7,66\ \text{m}}$
  \task $a=\frac{5}{0{,}25}=\res{20\ \text{cm}}$, \ $b=\sqrt{20^2-5^2}=\res{5\sqrt{15}\approx19,36\ \text{cm}}$, \ $\hat C\approx\res{14^\circ28'39''}$, \ $\hat B\approx\res{75^\circ31'21''}$
\end{tasks}

\ex{1} \begin{tasks}(2)
  \task $h=18\sin62^\circ\approx\res{15,89\ \text{cm}}$
  \task $h=28\sin35^\circ\approx\res{16,06\ \text{cm}}$
\end{tasks}

\ex{1} $\cos\alpha=\frac{1{,}2}{3}=0{,}4\Rightarrow\alpha\approx\res{66^\circ25'19''}$

\ex{1} $\tg\alpha=\frac{20}{25}=0{,}8\Rightarrow\alpha\approx\res{38^\circ39'35''}$

\ex{2} La altura lo divide en dos triángulos rectángulos de catetos 9 y 10: $\tg\hat B=\frac{10}{9}\Rightarrow$ ángulos iguales $\approx\res{48^\circ0'46''}$; \ ángulo desigual $\approx180^\circ-2\cdot48^\circ0'46''=\res{83^\circ58'28''}$.

\ex{2} Desnivel: $1265-785=480$ m. Pendiente: $\frac{480}{3000}=\res{0{,}16=16\,\%}$. \ Ángulo: $\tg^{-1}0{,}16\approx\res{9^\circ5'25''}$.

\ex{2} La distancia recorrida es la hipotenusa: $\sin\alpha=\frac{2\,160}{18\,410}\Rightarrow\alpha\approx\res{6^\circ44'16''}$.

\ex{2} El triángulo es isósceles: la altura divide el ángulo en dos de $25^\circ$. \ $r=2\cdot12\sin25^\circ\approx\res{10,14\ \text{cm}}$.

\ex{2} El ángulo recto está en $C$: $\overline{AC}=\frac{120}{\tg37^\circ}\approx159,25$ m, \ $\overline{AB}=\frac{120}{\sin37^\circ}\approx199,4$ m. Diferencia: $\approx\res{40,15\ \text{m}}$.


% =====================================================
\seccio{6}{Problemas de trigonometría}

\exf{2} $h=x\tg35^\circ=(x+200)\tg20^\circ\Rightarrow x=\dfrac{200\tg20^\circ}{\tg35^\circ-\tg20^\circ}\approx216,5$ m; \ $h\approx\res{151,59\ \text{m}}$.

\ex{2} Desde el punto más próximo el ángulo es mayor: $h=x\tg46^\circ=(x+100)\tg33^\circ\Rightarrow x\approx168,19$ m; \ $h\approx\res{174,16\ \text{m}}$.

\ex{2} $h=x\tg76^\circ=(x+480)\tg45^\circ\Rightarrow x\approx159,43$ m, \ $h\approx639,43$ m sobre los puntos de observación. Altitud de la cima: $1200+h\approx\res{1839,43\ \text{m}}$.

\ex{2} $h=x\tg50^\circ=(x+5)\tg40^\circ\Rightarrow x\approx11,9$ m; \ $h\approx\res{14,18\ \text{m}}$.

\ex{2} Las ciudades están a lados opuestos: $h=x\tg36^\circ=(75-x)\tg12^\circ\Rightarrow x\approx16,98$ km; \ $h\approx\res{12,33\ \text{km}}$. \ Distancias: $\frac{h}{\sin36^\circ}\approx\res{20,98\ \text{km}}$ de $A$ y $\frac{h}{\sin12^\circ}\approx\res{59,32\ \text{km}}$ de $B$.

\ex{2} $h=x\tg60^\circ=(140-x)\tg40^\circ\Rightarrow x\approx45,69$ cm; \ la lámpara queda a $h\approx\res{79,14\ \text{cm}}$ del techo. \ Cuerda: $\frac{h}{\sin60^\circ}+\frac{h}{\sin40^\circ}\approx\res{214,49\ \text{cm}}$.

\ex{2} $h=x\tg60^\circ=(126-x)\tg30^\circ\Rightarrow\sqrt3\,x=\frac{126-x}{\sqrt3}\Rightarrow x=31{,}5$ m. \ Altura: $\res{\frac{63\sqrt3}{2}\approx54,56\ \text{m}}$. \ Cables: $\overline{AB}=\frac{31{,}5}{\cos60^\circ}=\res{63\ \text{m}}$, \ $\overline{BC}=\frac{94{,}5}{\cos30^\circ}=\res{63\sqrt3\approx109,12\ \text{m}}$. \ $\widehat{ABC}=180^\circ-60^\circ-30^\circ=\res{90^\circ}$.

\ex{2} $10\cos45^\circ+10\cos30^\circ=5\sqrt2+5\sqrt3\approx\res{15,73\ \text{m}}$

\ex{2} $25\sin30^\circ+30\sin50^\circ\approx12{,}5+22,98\approx\res{35,48\ \text{m}}$ (el tramo horizontal no cuenta).

\ex{2} \begin{enumerate}[label=\alph*),leftmargin=*,itemsep=2pt]\item $\tg36^\circ=\frac{5}{\text{ap}}\Rightarrow\text{ap}=\frac{5}{\tg36^\circ}\approx\res{6,88\ \text{cm}}$\item Son 5 triángulos isósceles de lados 6 y 6 y ángulo $72^\circ$: altura $6\sin72^\circ$; \ $S=5\cdot\frac{6\cdot6\sin72^\circ}{2}\approx\res{85,6\ \text{cm}^2}$\end{enumerate}

\ex{2} \begin{enumerate}[label=\alph*),leftmargin=*,itemsep=2pt]\item $h=\frac{2{,}5}{\tg27{,}5^\circ}\approx\res{4,8\ \text{cm}}$; \ $S=\frac{5h}{2}\approx\res{12,01\ \text{cm}^2}$\item Lados iguales: $\frac{17}{\sin25^\circ}\approx40,23$ cm; \ $P\approx\res{114,45\ \text{cm}}$; \ $h=\frac{17}{\tg25^\circ}\approx36,46$ cm, \ $S\approx\res{619,76\ \text{cm}^2}$\end{enumerate}

\ex{2} \begin{enumerate}[label=\alph*),leftmargin=*,itemsep=2pt]\item $S=12\cdot12\sin43^\circ\approx\res{98,21\ \text{cm}^2}$\item Las diagonales se cortan perpendicularmente en el punto medio: $\tg\frac{\alpha}{2}=\frac64\Rightarrow\alpha\approx\res{112^\circ37'12''}$ y $\beta\approx\res{67^\circ22'48''}$; \ lado $\sqrt{6^2+4^2}=2\sqrt{13}$, \ $P=\res{8\sqrt{13}\approx28,84\ \text{cm}}$; \ $S=\frac{12\cdot8}{2}=\res{48\ \text{cm}^2}$\end{enumerate}

\ex{3} Trazando la altura sobre el lado de 15 cm: $h=10\sin35^\circ\approx5,74$, proyección $10\cos35^\circ\approx8,19$. \ Diagonal menor: $\sqrt{h^2+(15-8,19)^2}\approx\res{8,9\ \text{cm}}$. \ Área: $15\cdot h\approx\res{86,04\ \text{cm}^2}$.

\ex{3} Trazamos la altura desde $A$ sobre $BC$: $h=245\sin57^\circ\approx205,47$ m, $\overline{CH}=245\cos57^\circ\approx133,44$ m, $\overline{HB}\approx524,56$ m. \ $\overline{AB}=\sqrt{h^2+\overline{HB}^2}\approx\res{563,37\ \text{m}}$.

\ex{3} La portería va de 27,5 m a 32,5 m de altura sobre la línea de banda y $A$ está a 50 m de la línea de fondo. Ángulo de $AB$ con la línea de banda: $\tg^{-1}\frac{27{,}5}{50}\approx28^\circ48'39''$; \ ángulo de $AC$: $\tg^{-1}\frac{32{,}5}{50}\approx33^\circ1'26''$. \ $\alpha\approx\res{4^\circ12'47''}$.


% =====================================================
\seccio{7}{Razones trigonométricas de un ángulo cualquiera}

\exf{1} \begin{center}\renewcommand{\arraystretch}{1.35}\begin{tabular}{|c|C{2.2cm}|C{2.2cm}|C{2.2cm}|C{2cm}|}\hline $\alpha$ & $\sin\alpha$ & $\cos\alpha$ & $\tg\alpha$ & \textbf{Cuadrante}\\\hline $120^\circ$ & $\res{0,8660}$ & $\res{-0,5000}$ & $\res{-1,7321}$ & $\res{\text{II}}$\\\hline $220^\circ$ & $\res{-0,6428}$ & $\res{-0,7660}$ & $\res{0,8391}$ & $\res{\text{III}}$\\\hline $135^\circ$ & $\res{0,7071}$ & $\res{-0,7071}$ & $\res{-1,0000}$ & $\res{\text{II}}$\\\hline $300^\circ$ & $\res{-0,8660}$ & $\res{0,5000}$ & $\res{-1,7321}$ & $\res{\text{IV}}$\\\hline $240^\circ$ & $\res{-0,8660}$ & $\res{-0,5000}$ & $\res{1,7321}$ & $\res{\text{III}}$\\\hline $75^\circ$ & $\res{0,9659}$ & $\res{0,2588}$ & $\res{3,7321}$ & $\res{\text{I}}$\\\hline \end{tabular}\end{center}

\ex{1} \begin{tasks}(3)
  \task \rest{III o IV}
  \task \rest{II o III}
  \task \rest{II o IV}
\end{tasks}

\ex{1} \begin{tasks}(4)
  \task $\res{-1}$
  \task $\res{1}$
  \task $\res{1}$
  \task $\res{1}$
\end{tasks}

\ex{2} \begin{tasks}(1)
  \task $\sin\alpha=\sqrt{1-0{,}49}\approx\res{0,7141}$, \ $\tg\alpha\approx\res{1,0202}$
  \task $\frac{1}{\cos^2\beta}=9+1=10$, \ $\cos\beta=\res{-\frac{\sqrt{10}}{10}}$ (negativo en el II), \ $\sin\beta=\tg\beta\cdot\cos\beta=\res{\frac{3\sqrt{10}}{10}}$
  \task $\cos\gamma=-\sqrt{1-\frac{9}{25}}=\res{-\frac45}$ (negativo en el III), \ $\tg\gamma=\frac{-3/5}{-4/5}=\res{\frac34}$
  \task $\sin\delta=-\sqrt{1-\frac{25}{169}}=\res{-\frac{12}{13}}$ (negativo en el IV), \ $\tg\delta=\res{-\frac{12}{5}}$
\end{tasks}

\ex{2} $\cos^2\alpha=1-0{,}34^2=0{,}8844\Rightarrow\cos\alpha=\pm0,9404$. Como el seno es positivo, $\alpha$ puede ser del I o del II cuadrante: \ I: $\cos\alpha\approx\res{0,9404}$, $\tg\alpha\approx\res{0,3615}$ \ ($\alpha\approx19^\circ52'37''$); \ II: $\cos\alpha\approx\res{-0,9404}$, $\tg\alpha\approx\res{-0,3615}$ \ ($\alpha\approx160^\circ7'23''$).

\ex{2} \begin{tasks}(2)
  \task $\res{210^\circ}$, \ $\res{330^\circ}$
  \task $\res{135^\circ}$, \ $\res{315^\circ}$
  \task $\res{36^\circ52'12''}$, \ $\res{323^\circ7'48''}$
  \task $\res{64^\circ9'29''}$, \ $\res{115^\circ50'31''}$
\end{tasks}

\ex{2} Seno negativo y coseno positivo: \rest{cuarto cuadrante}. \ $\beta=360^\circ-30^\circ=\res{330^\circ}$.

\ex{1} \begin{tasks}(2)
  \task $1750^\circ=4\cdot360^\circ+310^\circ$: \ $\res{310^\circ}$
  \task $1920^\circ=5\cdot360^\circ+120^\circ$: \ $\res{120^\circ}$
  \task $-800^\circ=-3\cdot360^\circ+280^\circ$: \ $\res{280^\circ}$
  \task $3825^\circ=10\cdot360^\circ+225^\circ$: \ $\res{225^\circ}$
\end{tasks}


% =====================================================
\seccio{8}{Reducción al primer cuadrante}

\exf{1} \begin{tasks}(1)
  \task $135^\circ=180^\circ-45^\circ$: \ $\sin=\res{\frac{\sqrt2}{2}}$, $\cos=\res{-\frac{\sqrt2}{2}}$, $\tg=\res{-1}$
  \task $240^\circ=180^\circ+60^\circ$: \ $\sin=\res{-\frac{\sqrt3}{2}}$, $\cos=\res{-\frac12}$, $\tg=\res{\sqrt3}$
  \task $-30^\circ=-360^\circ+330^\circ$; \ $330^\circ=360^\circ-30^\circ$: \ $\sin=\res{-\frac12}$, $\cos=\res{\frac{\sqrt3}{2}}$, $\tg=\res{-\frac{\sqrt3}{3}}$
\end{tasks}

\ex{2} \begin{tasks}(1)
  \task $155^\circ=180^\circ-25^\circ$: \ $\sin155^\circ=\res{\sin25^\circ}$, $\cos155^\circ=\res{-\cos25^\circ}$, $\tg155^\circ=\res{-\tg25^\circ}$
  \task $240^\circ=180^\circ+60^\circ$: \ $\sin=\res{-\frac{\sqrt3}{2}}$, $\cos=\res{-\frac12}$, $\tg=\res{\sqrt3}$
  \task $390^\circ=360^\circ+30^\circ$: \ $\sin=\res{\frac12}$, $\cos=\res{\frac{\sqrt3}{2}}$, $\tg=\res{\frac{\sqrt3}{3}}$
  \task $225^\circ=180^\circ+45^\circ$: \ $\sin=\res{-\frac{\sqrt2}{2}}$, $\cos=\res{-\frac{\sqrt2}{2}}$, $\tg=\res{1}$
  \task $1920^\circ=5\cdot360^\circ+120^\circ$; \ $120^\circ=180^\circ-60^\circ$: \ $\sin=\res{\frac{\sqrt3}{2}}$, $\cos=\res{-\frac12}$, $\tg=\res{-\sqrt3}$
  \task $-30^\circ=-360^\circ+330^\circ$; \ $330^\circ=360^\circ-30^\circ$: \ $\sin=\res{-\frac12}$, $\cos=\res{\frac{\sqrt3}{2}}$, $\tg=\res{-\frac{\sqrt3}{3}}$
\end{tasks}

\ex{2} \begin{tasks}(1)
  \task $-40^\circ=-360^\circ+320^\circ$; \ $320^\circ=360^\circ-40^\circ$: \ $\sin(-40^\circ)=\res{-\sin40^\circ}$, $\cos(-40^\circ)=\res{\cos40^\circ}$, $\tg(-40^\circ)=\res{-\tg40^\circ}$
  \task $385^\circ=360^\circ+25^\circ$: \ $\sin385^\circ=\res{\sin25^\circ}$, $\cos385^\circ=\res{\cos25^\circ}$, $\tg385^\circ=\res{\tg25^\circ}$
  \task $-185^\circ=-360^\circ+175^\circ$; \ $175^\circ=180^\circ-5^\circ$: \ $\sin(-185^\circ)=\res{\sin5^\circ}$, $\cos(-185^\circ)=\res{-\cos5^\circ}$, $\tg(-185^\circ)=\res{-\tg5^\circ}$
  \task $835^\circ=2\cdot360^\circ+115^\circ$; \ $115^\circ=180^\circ-65^\circ$: \ $\sin835^\circ=\res{\sin65^\circ}$, $\cos835^\circ=\res{-\cos65^\circ}$, $\tg835^\circ=\res{-\tg65^\circ}$
  \task $934^\circ=2\cdot360^\circ+214^\circ$; \ $214^\circ=180^\circ+34^\circ$: \ $\sin934^\circ=\res{-\sin34^\circ}$, $\cos934^\circ=\res{-\cos34^\circ}$, $\tg934^\circ=\res{\tg34^\circ}$
\end{tasks}

\ex{2} \begin{tasks}(2)
  \task $145^\circ$: \ $\sin\approx\res{0,57}$, $\cos\approx\res{-0,82}$, $\tg\approx\res{-0,7}$
  \task $215^\circ$: \ $\sin\approx\res{-0,57}$, $\cos\approx\res{-0,82}$, $\tg\approx\res{0,7}$
  \task $325^\circ$: \ $\sin\approx\res{-0,57}$, $\cos\approx\res{0,82}$, $\tg\approx\res{-0,7}$
  \task $55^\circ$: \ $\sin\approx\res{0,82}$, $\cos\approx\res{0,57}$, $\tg\approx\res{1,43}$
  \task $-35^\circ$: \ $\sin\approx\res{-0,57}$, $\cos\approx\res{0,82}$, $\tg\approx\res{-0,7}$
  \task $395^\circ$: \ $\sin\approx\res{0,57}$, $\cos\approx\res{0,82}$, $\tg\approx\res{0,7}$
\end{tasks}

\ex{3} \begin{tasks}(1)
  \task $120^\circ=180^\circ-60^\circ$: \ $\sin120^\circ=\sin60^\circ=\res{\cos30^\circ}$, \ $\cos120^\circ=-\cos60^\circ=\res{-\sin30^\circ}$, \ $\tg120^\circ=-\tg60^\circ=\res{-\dfrac{1}{\tg30^\circ}}$
  \task $430^\circ=360^\circ+70^\circ$: \ $\sin430^\circ=\sin70^\circ=\res{\cos20^\circ}$, \ $\cos430^\circ=\res{\sin20^\circ}$, \ $\tg430^\circ=\res{\dfrac{1}{\tg20^\circ}}$
  \task $-800^\circ=-3\cdot360^\circ+280^\circ$ y $280^\circ=360^\circ-80^\circ$: \ $\sin(-800^\circ)=-\sin80^\circ=\res{-\cos10^\circ}$, \ $\cos(-800^\circ)=\cos80^\circ=\res{\sin10^\circ}$, \ $\tg(-800^\circ)=-\tg80^\circ=\res{-\dfrac{1}{\tg10^\circ}}$
  \task $\sin180^\circ=\res{\sin0^\circ=0}$, \ $\cos180^\circ=\res{-\cos0^\circ=-1}$, \ $\tg180^\circ=\res{\tg0^\circ=0}$
\end{tasks}

\ex{3} $\sin\alpha-\sin\alpha+\cos\alpha-\cos\alpha=\res{0}$


\end{document}
